\section{第5章\quad 方差缩减技术}

\subsection{故事：MC 的"省钱秘籍"}

蒙特卡洛估计的误差 $\propto\sigma/\sqrt n$。想把误差减小到 $1/10$？
样本量要乘 100——这是"花钱"的笨办法。
聪明的办法是\textbf{在 $n$ 不变的前提下把 $\sigma$ 压下去}：
同样的计算量换更高的精度，这就是方差缩减（variance reduction）。

本章的统一目标：估计积分/期望
\[
I=\int_0^1 f(x)\,\d x=\E\,f(U),\qquad U\sim U(0,1).
\]
每种技术都是一种"省钱"思路，学的时候抓住两点：\textbf{为什么方差变小了}、\textbf{怎么操作}。

\subsection{出发点：投点法 vs 平均值法}

\paragraph{投点法（hit-or-miss）。}
设 $0\le f\le1$。往单位正方形里均匀投 $n$ 个点 $(U_i,V_i)$，
数落在曲线下方（$V_i\le f(U_i)$）的比例：
\[
\hat I_1=\frac1n\sum_i I\big(V_i\le f(U_i)\big),\qquad
\Var(\hat I_1)=\frac{I(1-I)}{n}.
\]

\paragraph{平均值法（sample mean）。}
直接把 $I$ 看成期望，取
\[
\hat I_2=\frac1n\sum_i f(U_i),\qquad
\Var(\hat I_2)=\frac1n\Big(\int_0^1 f^2\,\d x-I^2\Big).
\]

\begin{kbox}{定理：平均值法永不吃亏}
若 $0\le f(x)\le1$，则 $\Var(\hat I_2)\le\Var(\hat I_1)$。

\textbf{证明}：$0\le f\le1\Rightarrow f^2\le f$，故
$\int f^2\,\d x\le\int f\,\d x=I$，于是
\[
\Var(\hat I_2)=\frac{\int f^2-I^2}{n}\le\frac{I-I^2}{n}=\frac{I(1-I)}{n}=\Var(\hat I_1).\qquad\blacksquare
\]
\textbf{直觉}：投点法把 $f(U_i)$ 这个连续信息压缩成"中/不中"一个比特，浪费了信息；
平均值法用全了，自然更准。
\end{kbox}

\begin{exbox}[两法方差对比]
估计 $I=\int_0^1 x\,\d x=\frac12$。分别算投点法与平均值法的方差，求方差缩减比。
\end{exbox}
\begin{sol}
投点法：$\Var(\hat I_1)=\frac{I(1-I)}{n}=\frac{1/4}{n}=\frac{0.25}{n}$。

平均值法：$\int_0^1x^2\d x=\frac13$，
$\Var(\hat I_2)=\frac1n\big(\frac13-\frac14\big)=\frac{1/12}{n}\approx\frac{0.083}{n}$。

缩减比 $=0.25/0.083=3$：平均值法等于白赚了 3 倍样本量。
\end{sol}

\subsection{对偶变量法（antithetic variates）}

\paragraph{直觉先行。}
两个估计量平均的方差：
\[
\Var\Big(\frac{X+Y}{2}\Big)=\frac{\Var X+\Var Y+2\Cov(X,Y)}{4}.
\]
若 $X,Y$ \textbf{负相关}，$\Cov<0$，平均后的方差比独立时还小——
"一个偏高时另一个偏低，互相抵消"。

\begin{ebox}{操作}
$U$ 与 $1-U$ 同为 $U(0,1)$ 但完全负相关。取
\[
\hat I=\frac1{2n}\sum_{i=1}^n\big[f(U_i)+f(1-U_i)\big].
\]
当 $f$ \textbf{单调}时，$f(U)$ 与 $f(1-U)$ 负相关，必有方差缩减；
且每对只花一个随机数，等于"买一送一"。
\end{ebox}

\begin{pbox}
抬桌子类比：一个人端盘子会抖（误差大），两个人抬且"你往左晃时我恰好往右晃"
（负相关），盘子就稳了。$U$ 和 $1-U$ 就是这样一对"反着抖"的搭档：
$U$ 抽到偏大的，$1-U$ 必然偏小，$f$ 单调时两个函数值一高一低，平均后误差大幅抵消。
\end{pbox}

\begin{exbox}[对偶法估计 $\int_0^1 e^x\d x$]
比较朴素平均值法与对偶法（同为 $2n$ 次函数求值）的方差。
已知 $\E e^U=e-1\approx1.718$，$\Var(e^U)=\frac{e^2-1}{2}-(e-1)^2\approx0.242$，
$\Cov(e^U,e^{1-U})=e-(e-1)^2\approx-0.234$。
\end{exbox}
\begin{sol}
朴素法用 $2n$ 个独立 $U$：$\Var=\frac{0.242}{2n}=\frac{0.121}{n}$。

对偶法：单对的方差
\[
\Var\Big(\frac{e^U+e^{1-U}}{2}\Big)=\frac{0.242+0.242+2\times(-0.234)}{4}
=\frac{0.016}{4}=0.0039 ,
\]
$n$ 对平均后 $\Var=\frac{0.0039}{n}$。
缩减比 $0.121/0.0039\approx31$ 倍！$e^x$ 单调性强，负相关抵消效果极好。
\end{sol}

\subsection{控制变量法（control variates）}

\paragraph{直觉先行。}
想估 $\E X$，手头另有一个与 $X$ 相关、且\textbf{期望已知}的变量 $Y$（$\E Y=\mu_Y$）。
$Y$ 的模拟误差 $Y-\mu_Y$ 是可观测的"风向标"——$X$ 与 $Y$ 正相关时，
$Y$ 偏高大概率 $X$ 也偏高，把这部分误差扣掉：
\[
X(c)=X-c\,(Y-\mu_Y).
\]
无论 $c$ 取多少都无偏（因 $\E(Y-\mu_Y)=0$），挑 $c$ 让方差最小。

\begin{kbox}{最优系数（必背推导）}
\[
\Var X(c)=\Var X-2c\Cov(X,Y)+c^2\Var Y
\]
是 $c$ 的开口向上抛物线，求导置零得
\[
c^*=\frac{\Cov(X,Y)}{\Var Y},\qquad
\Var X(c^*)=\Var X\,(1-\rho_{XY}^2).
\]
相关越强（$|\rho|\to1$），缩减越狠。实际中 $\Cov,\Var$ 用样本估计。
\end{kbox}

\begin{pbox}
考试生类比：想估某同学真实水平（$\E X$），但这次考题整体偏难。
怎么知道偏了多少？看一个"参照生"$Y$：他的真实水平 $\mu_Y$ 我们早就知道，
他这次考了 $Y$ 分，那么 $Y-\mu_Y$ 就是"这次考题的偏差"，
把它按比例 $c$ 从目标同学分数里扣掉即可。$c^*=\Cov/\Var$ 只是在问：
"两人受考题影响的联动有多强，就扣多大比例。"
\end{pbox}

\begin{exbox}[控制变量法估计 $\int_0^1 e^x\d x$]
取 $X=e^U$，控制变量 $Y=U$（$\E U=\frac12$ 已知）。
已知 $\Cov(e^U,U)=1-\frac{e-1}{2}\cdot1\approx0.1409$（即 $\E[Ue^U]-\E U\E e^U=1-(e-1)/2$），
$\Var U=\frac1{12}$，$\Var e^U\approx0.242$。求 $c^*$ 与方差缩减比。
\end{exbox}
\begin{sol}
\[
c^*=\frac{\Cov(e^U,U)}{\Var U}=\frac{0.1409}{1/12}\approx1.69 .
\]
相关系数 $\rho=\frac{0.1409}{\sqrt{0.242\times1/12}}\approx0.992$，
\[
\Var X(c^*)=0.242\times(1-0.992^2)\approx0.242\times0.016\approx0.0039 ,
\]
缩减约 60 倍。估计量为 $\frac1n\sum\big[e^{U_i}-1.69(U_i-\tfrac12)\big]$。
直线 $1.69(u-\frac12)+\text{const}$ 把 $e^u$ 拟合得很好，残差自然小。
\end{sol}

\subsection{条件期望法（Rao--Blackwell 化）}

\paragraph{直觉先行。}
全方差分解：
\[
\Var X=\E[\Var(X\mid Y)]+\Var[\E(X\mid Y)] .
\]
两项都非负，所以 $\Var[\E(X\mid Y)]\le\Var X$，而 $\E[\E(X\mid Y)]=\E X$（无偏不变）。
结论：\textbf{凡是能解析地算出来的部分就别模拟}——
用 $\E(X\mid Y)$ 替代 $X$，把"$Y$ 给定后 $X$ 的随机性"整个消灭掉。

\begin{pbox}
一句话：\textbf{能用笔算的部分就别抽签了}。估 $\pi$ 的例子里，
原来要抽两个随机数 $(U,V)$ 看点在不在圆内；但其实 $U$ 定下后，
"$V$ 落在圆内"的概率精确等于 $\sqrt{1-U^2}$，这一步根本不用抽——直接用公式。
少抽一次签，少一份随机性，方差自然小。
\end{pbox}

\begin{exbox}[条件期望法估计 $\pi$]
投点法估 $\pi$：$(U,V)$ 均匀于 $[0,1]^2$，$X=I(U^2+V^2\le1)$，$\E X=\pi/4$。
用条件期望法改进之。
\end{exbox}
\begin{sol}
对 $U$ 取条件：
\[
g(U)=\E[X\mid U]=\Pr(V\le\sqrt{1-U^2}\mid U)=\sqrt{1-U^2}.
\]
改用 $\hat\theta=\frac1n\sum\sqrt{1-U_i^2}$（只需一半随机数！）。
方差对比：$\Var X=\frac\pi4(1-\frac\pi4)\approx0.169$；
$\Var\sqrt{1-U^2}=\E(1-U^2)-(\pi/4)^2=\frac23-0.617\approx0.0498$。
缩减约 3.4 倍，且每次抽样成本减半。
\end{sol}

\subsection{重要抽样法（importance sampling）}

\paragraph{直觉先行。}
$I=\int f(x)\d x$ 中，若 $f$ 只在一小块区域大、其他地方几乎为 0，
均匀抽样的大部分点都"浪费"在不重要的区域。
不如按一个密度 $g$（形状接近 $f$）抽样，多去重要区域，再用权重纠偏：
\[
I=\int\frac{f(x)}{g(x)}\,g(x)\,\d x=\E_g\Big[\frac{f(X)}{g(X)}\Big]
\quad\Rightarrow\quad
\hat I=\frac1n\sum_i\frac{f(X_i)}{g(X_i)},\quad X_i\sim g .
\]

\begin{pbox}
找鱼类比：湖里的鱼（被积函数的"质量"）集中在东南角，均匀撒网是浪费；
聪明的做法是哪里鱼多往哪里多撒网（按 $g$ 抽样）。但多撒了网会高估鱼量，
所以每网收获要除以"这里撒网的加密倍数"（除以 $g$）来纠偏。
唯一的雷：有鱼的地方几乎不撒网（$f>0$ 处 $g\approx0$），一旦撞上权重爆炸。
\end{pbox}

\begin{kbox}{要点}
\begin{itemize}[leftmargin=2em]
  \item $g$ 与 $f/I$ 越像，$f/g$ 越接近常数，方差越小（理想极限：$g\propto|f|$ 时方差为 0，但需要知道 $I$，做不到，只能逼近）；
  \item \textbf{大忌}：$g$ 的尾巴比 $f$ 轻（$f>0$ 处 $g\approx0$），权重 $f/g$ 爆炸，方差反而无穷大。
\end{itemize}
\end{kbox}

\begin{exbox}[重要抽样估计 $\int_0^1 e^x\d x$]
取 $g(x)=\frac{2}{3}(1+x)$（$0\le x\le1$ 上的密度）。写出估计量并说明为何方差小。
\end{exbox}
\begin{sol}
$g$ 是归一化的 $1+x$：$\int_0^1(1+x)\d x=\frac32$，故 $g=\frac23(1+x)$。
抽 $X\sim g$（逆变换：$G(x)=\frac23(x+\frac{x^2}2)=u$，解 $x=\sqrt{1+3u}-1$），
\[
\hat I=\frac1n\sum_i\frac{e^{X_i}}{\frac23(1+X_i)}.
\]
因为 $e^x\approx1+x$（泰勒展开），比值 $\frac{e^x}{1+x}$ 在 $[0,1]$ 上仅在
$[1,\,e/2\approx1.36]$ 间变化，远比 $e^x\in[1,e]$ 平稳，方差显著小于朴素平均值法。
\end{sol}

\subsection{分层抽样法（stratified sampling）}

\paragraph{直觉先行。}
全国调查收入，纯随机抽样可能恰好多抽了富人——不如\textbf{先分层、层内定额}：
每层抽固定数量，层间差异（方差的"between"部分）被彻底消灭，只剩层内方差。

\begin{ebox}{操作（等分层为例）}
把 $[0,1]$ 切成 $m$ 段，第 $j$ 段抽 $n_j$ 个均匀点（$\sum n_j=n$），
\[
\hat I=\sum_{j=1}^m\frac{1}{m}\cdot\frac{1}{n_j}\sum_{i=1}^{n_j}f\big(X_{ij}\big),
\qquad X_{ij}\sim U\big(\tfrac{j-1}{m},\tfrac{j}{m}\big).
\]
由全方差分解，分层（按比例分配）后方差 $\le$ 简单随机抽样的方差；
按"层标准差大的多抽"（Neyman 分配）还能更优。
\end{ebox}

\begin{exbox}[分层抽样的效果]
估计 $I=\int_0^1x\,\d x$。把 $[0,1]$ 等分两层各抽 $n/2$ 个点，与不分层（$n$ 个点）比较方差。
\end{exbox}
\begin{sol}
不分层：$\Var=\frac{\Var U}{n}=\frac{1/12}{n}$。

分层：每层是宽 $\frac12$ 的均匀分布，层内方差 $\frac{(1/2)^2}{12}=\frac1{48}$。
\[
\Var(\hat I)=\frac14\Big(\frac{1/48}{n/2}+\frac{1/48}{n/2}\Big)=\frac{1}{48n}.
\]
缩减比 $\frac{1/12}{1/48}=4$ 倍。直觉：$f(x)=x$ 的差异主要在"层间"（左半小右半大），
分层把这部分方差直接清零。
\end{sol}

\subsection{公共随机数法（common random numbers）}

场景不同于前面：\textbf{比较}两个系统 $\theta_1-\theta_2$。
\[
\Var(X-Y)=\Var X+\Var Y-2\Cov(X,Y),
\]
让两个系统用\textbf{同一串随机数}模拟，制造\textbf{正}相关，差的方差就小——
"同一套考题考两个学生，分差更能反映真实水平"。
（注意与对偶变量法相反：比较问题要正相关，求和/平均问题要负相关。）

\subsection{本章考点地图}

\begin{ebox}{高频题型清单}
\begin{enumerate}[leftmargin=2em]
  \item \textbf{证明平均值法优于投点法}（$f^2\le f$ 三行证明，必考）；
  \item \textbf{对偶变量法}：写估计量+用 $\Cov<0$ 说明方差缩减；给数值算缩减比；
  \item \textbf{控制变量法}：推导 $c^*=\Cov/\Var$ 与 $\Var_{\min}=\Var X(1-\rho^2)$（高频大题）；
  \item \textbf{条件期望法}：全方差分解一行定理+构造 $\E(X\mid Y)$（如 $\sqrt{1-U^2}$ 估 $\pi$）；
  \item \textbf{重要抽样}：写加权估计量、说明 $g$ 的选取原则与尾部陷阱；
  \item \textbf{方法匹配概念题}：每种方法各自的适用场景与直觉一句话。
\end{enumerate}
\end{ebox}
